
%-------------------------------------------------------------------------
\chapter{INTRODUCTION}
%-------------------------------------------------------------------------

\section{Background}
The emergence of Large Language Models (LLMs) as practical tools for knowledge work has fundamentally altered expectations of what software can do. Systems like GPT-4, Claude, and Gemini demonstrated that a single model, trained on vast corpora of human knowledge, could reason, write, summarize, and converse at a level previously unimaginable \cite{anthropic2024mcp}. However, this capability came with a critical limitation: the model existed in isolation. It could only work with information placed directly into its context window. It could not browse the web, query a live database, read a file from disk, or call an API — not without a developer writing custom glue code to bridge the gap.

This limitation gave rise to the concept of \textbf{tool use} — the ability for an LLM to invoke external functions and receive their results. Early implementations were promising but deeply fragmented. Every AI provider had its own function-calling format. Every application had to build its own connectors. The ecosystem was a patchwork of incompatible, one-off integrations that were expensive to build and brittle to maintain \cite{anthropic2024mcp}.

The \textbf{Model Context Protocol (MCP)}, introduced by Anthropic in November 2024, is the industry's answer to this fragmentation. It is an open, vendor-neutral standard that defines a universal language for AI assistants to communicate with external tools and data sources. Just as USB standardized how peripherals connect to computers, MCP standardizes how AI models connect to the world \cite{anthropic2024mcp}.

\section{Problem Statement}
The core problem that MCP addresses is the \textbf{M×N Integration Problem}. In a world with M different AI models and N different tools or data sources, a naive approach requires building M×N custom integrations. Each integration is a bespoke engineering effort: a developer must understand the AI model's specific function-calling API, the tool's specific interface, and write code to translate between them. When the AI model updates its API, or the tool changes its schema, the integration breaks and must be rebuilt.

This creates a compounding maintenance burden. A company using three different AI assistants (Claude, GPT-4, Gemini) and ten internal tools (database, calendar, email, file system, CRM, ticketing system, code repository, analytics dashboard, document store, and search engine) would theoretically need to maintain 30 separate integrations. Each integration is a liability: a potential point of failure, a security surface, and a drain on engineering resources \cite{anthropic2024mcp}.

Beyond the engineering cost, this fragmentation limits what AI systems can actually accomplish. An AI assistant that cannot reliably access real-time data, execute actions, or maintain context across tool calls is fundamentally constrained in its usefulness for complex, multi-step tasks.

\section{Objectives}
The objectives of this report are as follows:
\begin{itemize}
    \item To trace the historical evolution of AI tool integration and understand why ad-hoc approaches failed.
    \item To provide a comprehensive understanding of the Model Context Protocol — its design philosophy, architecture, and core primitives.
    \item To examine the three-layer MCP architecture (Host, Client, Server) in depth, understanding the role and responsibilities of each component.
    \item To demonstrate the practical application of MCP through a fully worked example of building a custom MCP server.
    \item To evaluate MCP's significance as an infrastructure standard and its implications for the future of AI development.
\end{itemize}

\section{Scope}
This report covers the MCP specification as released by Anthropic in 2024, focusing on the conceptual architecture, the three core primitives (Tools, Resources, Prompts), the transport mechanisms (stdio and SSE), and the practical implementation of a server. It does not cover the full JSON-RPC 2.0 specification in exhaustive detail, nor does it benchmark MCP against alternative integration frameworks.

\section{Applications}
MCP has broad applicability across any domain where AI assistants need to interact with external systems:
\begin{itemize}
    \item \textbf{Software Development:} AI coding assistants that can read codebases, run tests, query documentation, and commit changes.
    \item \textbf{Enterprise Automation:} AI agents that can query CRM systems, update databases, send emails, and schedule meetings.
    \item \textbf{Research and Analysis:} AI systems that can search academic databases, retrieve papers, and synthesize findings from live sources.
    \item \textbf{Healthcare:} AI assistants that can securely query patient records, check drug interactions, and retrieve clinical guidelines.
    \item \textbf{Education:} Intelligent tutoring systems that can access course materials, grade submissions, and track student progress.
\end{itemize}

\section{Alignment with Sustainable Development Goals (SDGs)}
This work aligns with several United Nations Sustainable Development Goals:
\begin{itemize}
    \item \textbf{SDG 4: Quality Education:} By enabling AI assistants to reliably access educational resources and tools, MCP can power more capable and accurate intelligent tutoring systems.
    \item \textbf{SDG 9: Industry, Innovation, and Infrastructure:} MCP represents a foundational infrastructure innovation that reduces the cost of building AI-powered applications, democratizing access to capable AI tools for organizations of all sizes.
    \item \textbf{SDG 17: Partnerships for the Goals:} As an open standard, MCP fosters collaboration between AI providers, tool developers, and enterprises, creating a shared ecosystem that benefits all participants.
\end{itemize}

%-------------------------------------------------------------------------
\chapter{BEFORE MCP: THE AI INTEGRATION REVOLUTION AND ITS CHAOS}
%-------------------------------------------------------------------------

\section{The Rise of Tool-Using AI}
The story of MCP begins not with a protocol, but with a problem that grew slowly and then all at once. In the early days of LLMs, the models were purely generative: given text, produce text. Their value was in language understanding and generation, and their limitations were accepted as fundamental. They were oracles you could consult, not agents you could deploy \cite{anthropic2024mcp}.

This changed with the introduction of \textbf{function calling} by OpenAI in 2023. For the first time, a developer could describe a set of functions to the model — their names, parameters, and descriptions — and the model could decide to invoke one of them, returning a structured JSON object that the application could execute. This was a paradigm shift. The model was no longer just answering questions; it was directing actions.

The implications were immediately understood by the industry. An AI that could call functions could, in principle, do anything a computer could do: search the web, query databases, send messages, control software, and interact with APIs. The race to build "AI agents" — autonomous systems that could complete multi-step tasks — began in earnest.

\section{The M×N Integration Chaos}
As the ecosystem of AI models and tools exploded, a structural problem emerged with equal speed. Consider the landscape by late 2023: there were multiple competing AI providers (OpenAI, Anthropic, Google, Mistral, Cohere), each with their own function-calling API format. There were thousands of tools and services that developers wanted to connect to these models. And there was no standard.

\begin{figure}[H]
    \centering
    % IMAGE PLACEHOLDER 1
    % Suggested image: A diagram showing the M×N integration problem.
    % Draw a grid or web where on the left side you have AI Models (Claude, GPT-4, Gemini, Llama)
    % and on the right side you have Tools (GitHub, Slack, Database, File System, Calendar, CRM).
    % Draw arrows from EVERY model to EVERY tool, creating a tangled web of N*M connections.
    % Label it "The M×N Integration Problem: 4 Models × 6 Tools = 24 Custom Integrations"
    \includegraphics[width=0.85\textwidth]{mn_integration_chaos.png}
    \caption{The M×N Integration Problem: Every AI model requires a custom connector to every tool, resulting in an exponentially growing maintenance burden.}
    \label{fig:mn_chaos}
\end{figure}

Each integration was a custom engineering project. A developer building a GitHub integration for Claude had to learn Anthropic's tool-use API format. If they then wanted the same integration to work with GPT-4, they had to rewrite it for OpenAI's format. If they wanted it to work with Gemini, they rewrote it again. The same tool, three times, in three different formats, with three different maintenance burdens \cite{anthropic2024mcp}.

This is the M×N problem: with M models and N tools, you need up to M×N integrations. As M and N both grew rapidly, the total integration surface area grew quadratically. The industry was building on sand.

\section{The Anatomy of a Pre-MCP Integration}
To understand why this was so painful, it is worth examining what a pre-MCP integration actually looked like. Suppose a developer wanted to give an AI assistant the ability to query a PostgreSQL database. The process involved several distinct, brittle steps.

\textbf{Step 1: Define the tool schema.} The developer had to write a JSON schema describing the tool — its name, description, and parameters — in the specific format required by the AI provider. OpenAI's format used a \texttt{functions} array with a specific structure. Anthropic's format used a \texttt{tools} array with a slightly different structure. These were not interchangeable.

\textbf{Step 2: Handle the model's response.} When the model decided to call the tool, it returned a structured response. The developer had to parse this response, extract the function name and arguments, and route the call to the appropriate handler. Each provider's response format was different.

\textbf{Step 3: Execute the tool and return results.} The developer executed the database query and had to format the result in the specific way the AI provider expected for tool results. Again, each provider had its own format.

\textbf{Step 4: Manage the conversation loop.} The developer had to append the tool result to the conversation history and re-invoke the model, managing the multi-turn loop manually. This logic was entirely application-specific.

Every one of these steps was provider-specific. A change in the provider's API could break any of them. And this entire process had to be repeated for every tool, for every provider.

\section{The Proliferation of Partial Solutions}
The industry did not sit idle in the face of this chaos. Several partial solutions emerged, each addressing some aspects of the problem while introducing new ones.

\textbf{LangChain} became the dominant framework for building LLM applications. It provided abstractions over multiple AI providers and a library of pre-built tool integrations. However, LangChain was a Python library, not a protocol. It solved the problem for Python developers but did nothing for developers in other languages. It also introduced its own abstractions and complexity, and its rapid evolution made it difficult to maintain applications built on it \cite{anthropic2024mcp}.

\textbf{OpenAI's Plugins} (later GPTs with Actions) attempted to create a standard for connecting ChatGPT to external services via OpenAPI specifications. This was a step toward standardization, but it was proprietary to OpenAI's ecosystem. A plugin built for ChatGPT did not work with Claude or Gemini.

\textbf{Agent frameworks} like AutoGPT, BabyAGI, and later LangGraph and CrewAI provided higher-level abstractions for building autonomous agents. But they all faced the same underlying problem: the tool integration layer was still fragmented and provider-specific.

The fundamental issue was that all of these solutions were \textit{libraries} or \textit{frameworks}, not \textit{protocols}. A library is a piece of code you import into your application. A protocol is a standard that defines how two independent systems communicate. The industry needed a protocol.

\begin{table}[H]
    \centering
    \caption{Comparison of Pre-MCP Integration Approaches}
    \label{tab:pre_mcp_comparison}
    \begin{tabular}{|>{\raggedright\arraybackslash}p{2.8cm}|>{\raggedright\arraybackslash}p{3.0cm}|>{\raggedright\arraybackslash}p{3.5cm}|>{\raggedright\arraybackslash}p{3.5cm}|}
        \hline
        \textbf{Approach} & \textbf{Provider} & \textbf{Strengths} & \textbf{Limitations} \\ \hline
        Native Function Calling & OpenAI, Anthropic, Google & Direct, low-latency, full control & Provider-specific, no reuse across models \\ \hline
        LangChain & Open Source & Multi-provider, large tool library & Python-only, high abstraction overhead, rapid API churn \\ \hline
        OpenAI Plugins / GPT Actions & OpenAI & Standardized via OpenAPI & Proprietary to ChatGPT ecosystem \\ \hline
        Custom Agent Frameworks & Various & Flexible, composable & Each framework reinvents the wheel \\ \hline
        Direct API Calls & N/A & Maximum control & No AI awareness, manual context management \\ \hline
    \end{tabular}
\end{table}

\section{The Tipping Point}
By mid-2024, the AI industry had reached a tipping point. The promise of AI agents — systems that could autonomously complete complex, multi-step tasks — was clear. But the infrastructure to support them was a mess. The most capable AI models in the world were being held back not by their intelligence, but by the lack of a standard way to connect them to the world. The stage was set for MCP.

%-------------------------------------------------------------------------
\chapter{WHAT IS THE MODEL CONTEXT PROTOCOL?}
%-------------------------------------------------------------------------

\section{Definition and Core Philosophy}
The \textbf{Model Context Protocol (MCP)} is an open standard, published by Anthropic in November 2024, that defines a universal communication protocol between AI applications (hosts) and external capability providers (servers). It is, at its core, a specification for how an AI assistant should ask for things from the outside world, and how the outside world should respond \cite{anthropic2024mcp}.

The design philosophy of MCP rests on three pillars:

\textbf{1. Universality:} MCP is designed to be model-agnostic and vendor-neutral. A tool built as an MCP server works with any AI host that implements the MCP client specification, regardless of which underlying LLM powers that host. This is the USB principle applied to AI: one standard, universal compatibility.

\textbf{2. Separation of Concerns:} MCP cleanly separates the AI reasoning layer from the tool execution layer. The AI model does not need to know how a tool is implemented. The tool does not need to know which AI model is calling it. They communicate through a well-defined interface, and each can evolve independently.

\textbf{3. Security and Control:} MCP is designed with a clear trust hierarchy. The human user, operating through the Host application, is always in control. The AI model can request tool calls, but the Host decides whether to execute them. Sensitive operations can require explicit user approval. The protocol is designed to make it easy to build safe, auditable AI systems.

\section{The USB Analogy}
The most intuitive way to understand MCP is through the USB analogy. Before USB, connecting a peripheral to a computer was a nightmare. Keyboards used PS/2 connectors. Mice used serial ports. Printers used parallel ports. Scanners used SCSI. Every device required a different port, a different driver, and a different configuration process. Adding a new peripheral often required opening the computer case.

USB solved this by defining a single, universal standard. Any USB device — keyboard, mouse, printer, camera, hard drive — plugs into any USB port on any computer. The operating system handles the driver abstraction. The user just plugs it in and it works.

\begin{figure}[H]
    \centering
    % IMAGE PLACEHOLDER 2
    % Suggested image: A two-panel comparison diagram.
    % LEFT PANEL: "Before MCP" - Show an AI model in the center with messy, different-colored
    % arrows going to different tools, each labeled with a different format (JSON-RPC, REST, GraphQL, etc.)
    % RIGHT PANEL: "After MCP" - Show an AI model in the center with clean, uniform arrows
    % going through a single "MCP" layer to all tools. Clean and organized.
    % Title: "MCP: The USB Standard for AI Tools"
    \includegraphics[width=0.9\textwidth]{mcp_usb_analogy.png}
    \caption{MCP as the USB standard for AI: replacing fragmented, incompatible integrations with a single universal protocol.}
    \label{fig:usb_analogy}
\end{figure}

MCP is USB for AI tools. An MCP server — whether it exposes a database, a file system, a web search API, or a custom business tool — plugs into any MCP-compatible AI host. The host handles the protocol abstraction. The developer just builds the server once and it works everywhere.

\section{The Three Core Primitives}
MCP defines three fundamental primitives that servers can expose to clients. These primitives cover the full range of ways an AI assistant might need to interact with the external world \cite{anthropic2024mcp}.

\subsection{Tools}
Tools are executable functions that the AI model can invoke to perform actions or retrieve computed results. They are the most powerful primitive and the closest analogue to function calling in traditional AI APIs. A tool has a name, a description (which the AI uses to decide when to call it), and a JSON Schema defining its input parameters.

Examples of tools include: \texttt{search\_web(query: string)}, \texttt{execute\_sql(query: string, database: string)}, \texttt{send\_email(to: string, subject: string, body: string)}, \texttt{create\_github\_issue(repo: string, title: string, body: string)}.

The critical distinction from raw function calling is that tools are defined on the \textit{server}, not in the \textit{application code}. This means the same tool definition can be used by any MCP-compatible host without modification.

\subsection{Resources}
Resources are data sources that the AI model can read to gain context. Unlike tools, resources are not actions — they do not change state. They are read-only data endpoints that the AI can use to ground its responses in real information. Resources are identified by URIs and can represent files, database records, API responses, or any other structured data.

Examples of resources include: \texttt{file:///project/README.md}, \texttt{db://customers/\{id\}}, \texttt{github://repos/owner/repo/issues}, \texttt{calendar://events/today}.

Resources solve a fundamental problem: how does an AI assistant get access to the specific context it needs for a task without the user having to manually copy-paste information into the chat? With resources, the AI can proactively fetch the data it needs.

\subsection{Prompts}
Prompts are reusable, parameterized prompt templates that servers can expose. They allow server developers to encode best practices for interacting with their service directly into the protocol. A prompt template might define the optimal way to ask the AI to analyze a database schema, review a pull request, or summarize a document in a specific format.

Examples of prompts include: \texttt{analyze\_database\_schema(schema: string)}, \texttt{review\_pull\_request(diff: string, context: string)}, \texttt{summarize\_document(document: string, format: "bullet\_points" | "paragraph")}.

Prompts are a subtle but powerful feature. They shift the expertise of how to use a tool from the user to the tool developer, making AI integrations more accessible and consistent.

\section{The Transport Layer}
MCP is transport-agnostic by design, but defines two primary transport mechanisms \cite{anthropic2024mcp}:

\textbf{stdio (Standard Input/Output):} The client launches the server as a subprocess and communicates via stdin/stdout. This is the simplest transport, ideal for local tools running on the same machine as the host. It requires no network configuration and has minimal overhead.

\textbf{SSE (Server-Sent Events) over HTTP:} The server runs as an HTTP server, and the client connects via HTTP. This transport supports remote servers and is suitable for cloud-hosted tools. It uses SSE for server-to-client streaming and HTTP POST for client-to-server messages.

Both transports use \textbf{JSON-RPC 2.0} as the message format — a lightweight, well-established standard for remote procedure calls over JSON.

%-------------------------------------------------------------------------
\chapter{HOW MCP SOLVES THE INTEGRATION PROBLEM}
%-------------------------------------------------------------------------

\section{From M×N to M+N}
The most important thing MCP does is transform the M×N integration problem into an M+N problem. Instead of needing M×N custom integrations, you need M MCP client implementations (one per AI host) and N MCP server implementations (one per tool). Once both sides implement the standard, any client can talk to any server.

\begin{figure}[H]
    \centering
    % IMAGE PLACEHOLDER 3
    % Suggested image: A mathematical/visual diagram showing the transformation.
    % LEFT: Show M=4 models and N=6 tools with M*N=24 arrows between them (messy web).
    % RIGHT: Show M=4 models connecting to a central "MCP Protocol" hub, which connects to N=6 tools.
    % Only M+N=10 connections total. Label clearly: "Before: M×N = 24 integrations" vs "After: M+N = 10 implementations"
    \includegraphics[width=0.9\textwidth]{mn_to_m_plus_n.png}
    \caption{MCP transforms the M×N integration problem into M+N: each AI host implements MCP once, and each tool implements MCP once. Any combination then works automatically.}
    \label{fig:mn_to_mplus_n}
\end{figure}

This is not merely a quantitative improvement — it is a qualitative shift. In the M×N world, adding a new AI model requires building N new integrations. In the M+N world, adding a new AI model requires building exactly one MCP client implementation. Every existing MCP server immediately becomes available to the new model. Similarly, adding a new tool requires building exactly one MCP server. Every existing MCP host immediately gains access to the new tool \cite{anthropic2024mcp}.

\section{Decoupling Development Lifecycles}
One of the most practically significant benefits of MCP is that it decouples the development lifecycles of AI hosts and tool servers. In the pre-MCP world, when Anthropic updated Claude's tool-use API, every integration built for Claude had to be updated. This created a coordination problem: tool developers had to track changes across multiple AI providers and update their integrations accordingly.

With MCP, the protocol is the stable interface. Anthropic can update Claude's internal architecture, improve its reasoning, or change its pricing — none of this affects MCP servers. The MCP specification is versioned and stable. Tool developers build to the spec, not to a specific model's API.

\section{Enabling the Ecosystem}
Perhaps the most transformative aspect of MCP is what it enables at the ecosystem level. Because MCP is an open standard, anyone can build an MCP server for any tool or service. This has already happened: within months of MCP's release, the community had built MCP servers for GitHub, Slack, Google Drive, PostgreSQL, Brave Search, Puppeteer (web automation), filesystem access, and dozens of other services \cite{anthropic2024mcp}.

This ecosystem effect is self-reinforcing. As more MCP servers are built, MCP-compatible AI hosts become more capable. As more AI hosts adopt MCP, the incentive to build MCP servers increases. The protocol creates a positive feedback loop that benefits the entire AI ecosystem.

\section{Security and Trust Model}
MCP's architecture enforces a clear security model. The Host application is the trust boundary. It controls which MCP servers the client connects to, which tools are exposed to the AI model, and whether tool calls require user confirmation. The AI model can request tool calls, but it cannot execute them directly — the Host mediates every interaction.

This design makes it straightforward to build safe AI systems. A Host can implement policies like: "always ask the user before executing any tool that modifies data," "never expose tools that access sensitive files to untrusted AI models," or "log all tool calls for audit purposes." These policies are enforced at the Host level, independent of the AI model's behavior.

\section{Standardizing Context Management}
Beyond tool execution, MCP standardizes how context is managed across a session. The protocol defines how servers can push updates to clients (notifications), how clients can subscribe to resource changes, and how the conversation history is maintained. This makes it possible to build AI assistants that maintain coherent, stateful interactions with external systems over extended sessions — a capability that was extremely difficult to implement reliably in the pre-MCP world.

%-------------------------------------------------------------------------
\chapter{MCP ARCHITECTURE}
%-------------------------------------------------------------------------

\section{Overview of the Three-Layer Architecture}
MCP defines a clean three-layer architecture: the \textbf{Host}, the \textbf{Client}, and the \textbf{Server}. Each layer has distinct responsibilities and communicates with adjacent layers through well-defined interfaces. Understanding this architecture is essential to understanding how MCP works in practice \cite{anthropic2024mcp}.

\begin{figure}[H]
    \centering
    % IMAGE PLACEHOLDER 4
    % Suggested image: A layered architecture diagram.
    % Draw three horizontal layers or three boxes arranged left to right:
    % [HOST (Claude Desktop / VS Code / Custom App)]
    %   contains → [MCP CLIENT 1] [MCP CLIENT 2] [MCP CLIENT 3]
    %   each client connects via "MCP Protocol (JSON-RPC 2.0)" to →
    % [MCP SERVER A (GitHub)] [MCP SERVER B (Database)] [MCP SERVER C (File System)]
    % Show the LLM inside the Host box. Show arrows for "Tool Request" and "Tool Result".
    % This is the most important diagram in the report.
    \includegraphics[width=\textwidth]{mcp_architecture_overview.png}
    \caption{The MCP three-layer architecture: a Host application manages multiple MCP Clients, each connected to a dedicated MCP Server via the standardized JSON-RPC 2.0 protocol.}
    \label{fig:mcp_architecture}
\end{figure}

The key insight of this architecture is the \textbf{one-to-many relationship}: a single Host can manage multiple Clients simultaneously, and each Client maintains a dedicated connection to one Server. This allows an AI assistant to have access to many different tools and data sources at the same time, all managed through a single, coherent interface.

\section{The Message Flow}
Before examining each layer in detail, it is useful to understand the end-to-end message flow for a typical tool call \cite{anthropic2024mcp}:

\begin{enumerate}
    \item The user sends a message to the Host application (e.g., "What are the open issues in my GitHub repository?").
    \item The Host passes the message to the LLM, along with a list of available tools (gathered from all connected MCP Servers).
    \item The LLM determines that it needs to call the \texttt{list\_issues} tool and returns a tool call request.
    \item The Host receives the tool call request and routes it to the appropriate MCP Client.
    \item The MCP Client sends a \texttt{tools/call} JSON-RPC request to the MCP Server.
    \item The MCP Server executes the tool (calls the GitHub API) and returns the result.
    \item The MCP Client passes the result back to the Host.
    \item The Host appends the tool result to the conversation and re-invokes the LLM.
    \item The LLM generates a final response based on the tool result.
    \item The Host displays the response to the user.
\end{enumerate}

This entire flow is mediated by the MCP protocol, making it consistent and predictable regardless of which specific tools or AI models are involved.

%-------------------------------------------------------------------------
\chapter{THE HOST}
%-------------------------------------------------------------------------

\section{Definition and Role}
The \textbf{Host} is the top-level application that the user interacts with directly. It is the orchestrator of the entire MCP system. The Host is responsible for managing the user interface, maintaining the conversation with the LLM, managing the lifecycle of MCP Clients, and enforcing the security and trust policies of the system \cite{anthropic2024mcp}.

Examples of Host applications include:
\begin{itemize}
    \item \textbf{Claude Desktop:} Anthropic's desktop application for Claude, which was the first major MCP Host.
    \item \textbf{VS Code with Copilot:} Microsoft's code editor with AI assistance, which adopted MCP for tool integration.
    \item \textbf{Cursor:} An AI-first code editor that uses MCP to connect to development tools.
    \item \textbf{Custom enterprise applications:} Any application built by a developer that embeds an LLM and implements the MCP Host specification.
\end{itemize}

\section{Responsibilities of the Host}
The Host has several critical responsibilities that distinguish it from a simple chat interface \cite{anthropic2024mcp}:

\textbf{LLM Management:} The Host is responsible for all communication with the underlying LLM. It constructs the prompts, manages the conversation history, and interprets the model's responses. When the model requests a tool call, the Host is responsible for recognizing this and routing the request appropriately.

\textbf{Client Lifecycle Management:} The Host creates and manages MCP Client instances. It decides which servers to connect to (typically based on user configuration), initiates the connection process, and handles disconnections and reconnections. A Host might manage a pool of persistent client connections or create clients on demand.

\textbf{Tool Aggregation:} The Host queries all connected MCP Clients for their available tools, resources, and prompts, and aggregates this information into a unified capability list that is presented to the LLM. The LLM sees a single, flat list of available tools, without needing to know which server each tool comes from.

\textbf{Security Enforcement:} The Host is the security boundary of the system. It decides which tools to expose to the LLM, whether to require user confirmation for specific tool calls, and how to handle errors and unexpected behavior from servers. The Host can implement fine-grained access control policies.

\textbf{User Interface:} The Host presents the conversation to the user, displays tool call results, and provides mechanisms for the user to configure which MCP servers are connected.

\section{The Host as Trust Anchor}
A critical design principle of MCP is that the Host is the \textbf{trust anchor} of the system. The LLM is not trusted to make autonomous decisions about which tools to call or what data to access. Every tool call is mediated by the Host, which can apply policies, require user confirmation, and log all activity.

This design is intentional. LLMs can be manipulated through prompt injection attacks — malicious content in tool results that attempts to hijack the model's behavior. By making the Host the trust anchor, MCP ensures that even if the LLM is manipulated, the Host can detect and block unauthorized actions.

%-------------------------------------------------------------------------
\chapter{THE CLIENT}
%-------------------------------------------------------------------------

\section{Definition and Role}
The \textbf{MCP Client} is a lightweight component that lives inside the Host application and manages the connection to a single MCP Server. Each Client is responsible for the protocol-level communication with its server: sending requests, receiving responses, and handling the JSON-RPC message framing \cite{anthropic2024mcp}.

The Client is intentionally thin. It does not contain business logic, does not make decisions about when to call tools, and does not interpret tool results. Its sole responsibility is reliable, protocol-compliant communication with its server.

\section{Client Responsibilities}
\textbf{Connection Management:} The Client establishes and maintains the connection to the server. For stdio transport, this means launching the server process and managing its stdin/stdout streams. For SSE transport, this means establishing the HTTP connection and managing the event stream.

\textbf{Capability Discovery:} During the initialization handshake, the Client queries the server for its capabilities — the list of tools, resources, and prompts it exposes. This information is passed to the Host for aggregation.

\textbf{Request Routing:} When the Host needs to call a tool on a specific server, it routes the request to the appropriate Client. The Client formats the request as a JSON-RPC message and sends it to the server.

\textbf{Response Handling:} The Client receives responses from the server, validates them against the MCP specification, and passes them back to the Host. It handles error responses, timeouts, and unexpected disconnections gracefully.

\textbf{Notification Handling:} Servers can send unsolicited notifications to clients (e.g., "a resource has been updated"). The Client receives these notifications and forwards them to the Host, which can use them to update the AI's context.

\section{The Initialization Handshake}
When a Client connects to a Server, they perform a structured initialization handshake \cite{anthropic2024mcp}:

\begin{enumerate}
    \item The Client sends an \texttt{initialize} request, declaring its protocol version and capabilities.
    \item The Server responds with its protocol version, capabilities, and server information.
    \item The Client sends an \texttt{initialized} notification to confirm the handshake is complete.
    \item The Client can now send \texttt{tools/list}, \texttt{resources/list}, and \texttt{prompts/list} requests to discover the server's offerings.
\end{enumerate}

This handshake ensures that both sides agree on the protocol version and understand each other's capabilities before any tool calls are made.

\begin{figure}[H]
    \centering
    % IMAGE PLACEHOLDER 5
    % Suggested image: A sequence diagram (UML-style) showing the initialization handshake.
    % Three columns: HOST | CLIENT | SERVER
    % Arrows showing:
    % HOST → CLIENT: "Create Client, connect to server"
    % CLIENT → SERVER: "initialize {protocolVersion, capabilities}"
    % SERVER → CLIENT: "initialize response {protocolVersion, capabilities, serverInfo}"
    % CLIENT → SERVER: "initialized (notification)"
    % CLIENT → SERVER: "tools/list"
    % SERVER → CLIENT: "tools/list response {tools: [...]}"
    % CLIENT → HOST: "Available tools: [list_issues, create_issue, ...]"
    \includegraphics[width=0.85\textwidth]{mcp_handshake_sequence.png}
    \caption{The MCP initialization handshake sequence: Client and Server negotiate protocol version and exchange capability information before any tool calls are made.}
    \label{fig:handshake}
\end{figure}

%-------------------------------------------------------------------------
\chapter{THE SERVER}
%-------------------------------------------------------------------------

\section{Definition and Role}
The \textbf{MCP Server} is the component that exposes capabilities to the AI ecosystem. It is the bridge between the MCP protocol and the actual tools, data sources, and services that the AI needs to access. A server can be a simple script that wraps a single API, or a complex service that exposes dozens of tools and resources \cite{anthropic2024mcp}.

The Server is where the real work happens. It contains the business logic for executing tools, fetching resources, and rendering prompts. It is responsible for authentication with external services, error handling, and data transformation.

\section{Server Responsibilities}
\textbf{Tool Implementation:} The Server implements the actual logic for each tool it exposes. When it receives a \texttt{tools/call} request, it executes the corresponding function, handles any errors, and returns the result in the MCP-specified format.

\textbf{Resource Serving:} The Server implements handlers for each resource URI it exposes. When it receives a \texttt{resources/read} request, it fetches the requested data (from a file, database, API, etc.) and returns it as structured content.

\textbf{Prompt Rendering:} The Server implements prompt templates and renders them with the provided parameters when a \texttt{prompts/get} request is received.

\textbf{Authentication:} The Server handles authentication with external services. It may use API keys, OAuth tokens, or other credentials to authenticate with the services it wraps. These credentials are managed by the server and are never exposed to the AI model.

\textbf{Error Handling:} The Server is responsible for graceful error handling. If a tool call fails (e.g., the external API is unavailable), the server returns a structured error response that the Client can pass back to the Host.

\section{Server Types}
MCP servers can be categorized by their deployment model \cite{anthropic2024mcp}:

\textbf{Local Servers (stdio):} These servers run as subprocesses on the same machine as the Host. They are launched by the Host when needed and communicate via stdin/stdout. Local servers are ideal for tools that access local resources (file system, local databases, local development tools). They have no network overhead and can access local credentials securely.

\textbf{Remote Servers (SSE):} These servers run as HTTP services, potentially on remote machines or in the cloud. They are suitable for tools that wrap cloud APIs or need to be shared across multiple users or machines. Remote servers require network connectivity and must handle authentication carefully.

\begin{figure}[H]
    \centering
    % IMAGE PLACEHOLDER 6
    % Suggested image: A diagram showing two types of MCP servers.
    % LEFT SIDE: "Local Server (stdio)" - Show a laptop with HOST and SERVER both inside it.
    % Arrow between them labeled "stdin/stdout". Server connects to "Local File System", "Local DB".
    % RIGHT SIDE: "Remote Server (SSE)" - Show a laptop with HOST, connected via "HTTP/SSE" over
    % the internet to a cloud server. Cloud server connects to "GitHub API", "Slack API", "Cloud DB".
    \includegraphics[width=0.9\textwidth]{mcp_server_types.png}
    \caption{Two deployment models for MCP Servers: local servers communicate via stdio for low-latency access to local resources, while remote servers use HTTP/SSE for cloud-hosted tools.}
    \label{fig:server_types}
\end{figure}

\section{The Tool Definition Schema}
Every tool exposed by an MCP server must be defined with a precise schema. This schema is what the AI model uses to understand what the tool does and how to call it. A well-written tool schema is critical for reliable AI behavior \cite{anthropic2024mcp}.

A tool definition consists of:
\begin{itemize}
    \item \textbf{name:} A unique identifier for the tool (e.g., \texttt{search\_issues}).
    \item \textbf{description:} A natural language description of what the tool does. This is read by the AI model to decide when to use the tool. It should be clear, specific, and include examples of when to use it.
    \item \textbf{inputSchema:} A JSON Schema object defining the tool's parameters — their names, types, descriptions, and which are required.
\end{itemize}

The quality of the description is arguably the most important factor in how well an AI model uses a tool. A vague description leads to incorrect tool selection. A precise description leads to reliable, accurate tool use.

%-------------------------------------------------------------------------
\chapter{BUILDING AN MCP SERVER: A PRACTICAL EXAMPLE}
%-------------------------------------------------------------------------

\section{Overview of the Example}
To make the MCP architecture concrete, this chapter walks through the complete implementation of a custom MCP server. The example server, called \textbf{Weather MCP Server}, exposes two tools: one that retrieves the current weather for a given city, and one that retrieves a 5-day forecast. This example is intentionally simple to focus on the MCP mechanics rather than the complexity of the underlying service \cite{anthropic2024mcp}.

The server is implemented in Python using the official \texttt{mcp} SDK provided by Anthropic. It uses the stdio transport, making it suitable for local use with Claude Desktop or any other MCP-compatible host.

\section{Project Setup}
The first step is to set up the Python project and install the MCP SDK.

\begin{lstlisting}[language=bash, caption={Setting up the MCP server project}]
# Create project directory
mkdir weather-mcp-server
cd weather-mcp-server

# Create virtual environment
python -m venv venv
source venv/bin/activate  # Windows: venv\Scripts\activate

# Install the MCP SDK and HTTP library
pip install mcp httpx

# Create the server file
touch server.py
\end{lstlisting}

\section{Implementing the Server}
The complete server implementation is shown below. Each section is explained in detail following the code.

\begin{lstlisting}[language=Python, caption={Complete Weather MCP Server implementation (server.py)}]
import asyncio
import httpx
from mcp.server import Server
from mcp.server.stdio import stdio_server
from mcp import types

# Initialize the MCP Server with a name
app = Server("weather-server")

# The base URL for the Open-Meteo API (free, no API key required)
GEOCODING_URL = "https://geocoding-api.open-meteo.com/v1/search"
WEATHER_URL = "https://api.open-meteo.com/v1/forecast"

async def get_coordinates(city: str) -> tuple[float, float]:
    """Helper: resolve a city name to lat/lon coordinates."""
    async with httpx.AsyncClient() as client:
        response = await client.get(
            GEOCODING_URL,
            params={"name": city, "count": 1, "language": "en"}
        )
        data = response.json()
        if not data.get("results"):
            raise ValueError(f"City '{city}' not found.")
        result = data["results"][0]
        return result["latitude"], result["longitude"]

# --- Tool 1: Get Current Weather ---
@app.list_tools()
async def list_tools() -> list[types.Tool]:
    """Declare all tools this server exposes."""
    return [
        types.Tool(
            name="get_current_weather",
            description=(
                "Retrieves the current weather conditions for a given city. "
                "Use this when the user asks about the current weather, "
                "temperature, or conditions in a specific location."
            ),
            inputSchema={
                "type": "object",
                "properties": {
                    "city": {
                        "type": "string",
                        "description": "The name of the city (e.g., 'London', 'Tokyo')"
                    }
                },
                "required": ["city"]
            }
        ),
        types.Tool(
            name="get_weather_forecast",
            description=(
                "Retrieves a 5-day weather forecast for a given city. "
                "Use this when the user asks about future weather, "
                "upcoming conditions, or a weather forecast."
            ),
            inputSchema={
                "type": "object",
                "properties": {
                    "city": {
                        "type": "string",
                        "description": "The name of the city"
                    },
                    "days": {
                        "type": "integer",
                        "description": "Number of forecast days (1-7, default 5)",
                        "default": 5
                    }
                },
                "required": ["city"]
            }
        )
    ]

@app.call_tool()
async def call_tool(
    name: str,
    arguments: dict
) -> list[types.TextContent]:
    """Handle tool execution requests from the MCP client."""

    if name == "get_current_weather":
        city = arguments["city"]
        lat, lon = await get_coordinates(city)

        async with httpx.AsyncClient() as client:
            response = await client.get(
                WEATHER_URL,
                params={
                    "latitude": lat,
                    "longitude": lon,
                    "current": "temperature_2m,wind_speed_10m,weathercode",
                    "temperature_unit": "celsius"
                }
            )
            data = response.json()
            current = data["current"]

        result = (
            f"Current weather in {city}:\n"
            f"  Temperature: {current['temperature_2m']}°C\n"
            f"  Wind Speed: {current['wind_speed_10m']} km/h\n"
            f"  Weather Code: {current['weathercode']}"
        )
        return [types.TextContent(type="text", text=result)]

    elif name == "get_weather_forecast":
        city = arguments["city"]
        days = arguments.get("days", 5)
        lat, lon = await get_coordinates(city)

        async with httpx.AsyncClient() as client:
            response = await client.get(
                WEATHER_URL,
                params={
                    "latitude": lat,
                    "longitude": lon,
                    "daily": "temperature_2m_max,temperature_2m_min,weathercode",
                    "forecast_days": days,
                    "temperature_unit": "celsius"
                }
            )
            data = response.json()
            daily = data["daily"]

        lines = [f"{days}-day forecast for {city}:"]
        for i in range(len(daily["time"])):
            lines.append(
                f"  {daily['time'][i]}: "
                f"High {daily['temperature_2m_max'][i]}°C, "
                f"Low {daily['temperature_2m_min'][i]}°C"
            )
        return [types.TextContent(type="text", text="\n".join(lines))]

    else:
        raise ValueError(f"Unknown tool: {name}")

# Entry point: run the server using stdio transport
async def main():
    async with stdio_server() as (read_stream, write_stream):
        await app.run(
            read_stream,
            write_stream,
            app.create_initialization_options()
        )

if __name__ == "__main__":
    asyncio.run(main())
\end{lstlisting}

\section{Code Walkthrough}
\subsection{Server Initialization}
The server is initialized with \texttt{Server("weather-server")}. The name is used for identification during the MCP handshake. The \texttt{@app.list\_tools()} decorator registers the function that returns the server's tool definitions. This function is called by the Client during capability discovery.

\subsection{Tool Definitions}
The \texttt{list\_tools} function returns a list of \texttt{types.Tool} objects. Each tool has a \texttt{name}, a \texttt{description}, and an \texttt{inputSchema}. Notice how the descriptions are written in natural language, explaining not just what the tool does but \textit{when} the AI should use it. This is critical for reliable tool selection.

\subsection{Tool Execution}
The \texttt{@app.call\_tool()} decorator registers the function that handles tool execution. When the Client sends a \texttt{tools/call} request, this function is invoked with the tool name and arguments. The function dispatches to the appropriate implementation based on the tool name.

\subsection{Transport and Entry Point}
The \texttt{stdio\_server()} context manager sets up the stdio transport. The \texttt{app.run()} call starts the server's main loop, which reads JSON-RPC messages from stdin and writes responses to stdout.

\section{Connecting to Claude Desktop}
To use this server with Claude Desktop, add the following configuration to the Claude Desktop configuration file (\texttt{claude\_desktop\_config.json}) \cite{anthropic2024mcp}:

\begin{lstlisting}[language=json, caption={Claude Desktop configuration for the Weather MCP Server}]
{
  "mcpServers": {
    "weather": {
      "command": "python",
      "args": ["/path/to/weather-mcp-server/server.py"],
      "env": {}
    }
  }
}
\end{lstlisting}

After restarting Claude Desktop, the weather tools will be available. The user can now ask Claude questions like "What's the weather in Tokyo?" or "Give me a 7-day forecast for London," and Claude will automatically invoke the appropriate tool.

\begin{figure}[H]
    \centering
    % IMAGE PLACEHOLDER 7
    % Suggested image: A screenshot or mockup of Claude Desktop showing the weather tool in action.
    % Show a chat interface where:
    % User: "What's the weather like in Tokyo right now?"
    % Claude: [Tool call indicator: "Using weather-server > get_current_weather"]
    % Claude: "The current weather in Tokyo is 22°C with wind speeds of 15 km/h..."
    % This shows the end-to-end user experience of MCP in action.
    \includegraphics[width=0.85\textwidth]{claude_desktop_weather.png}
    \caption{Claude Desktop using the Weather MCP Server: the AI automatically invokes the \texttt{get\_current\_weather} tool when the user asks about weather conditions.}
    \label{fig:claude_desktop}
\end{figure}

\section{Testing the Server}
Before connecting to a Host, the server can be tested directly using the MCP Inspector, a developer tool provided by Anthropic:

\begin{lstlisting}[language=bash, caption={Testing the server with MCP Inspector}]
# Install the MCP Inspector
npx @modelcontextprotocol/inspector python server.py

# The Inspector opens a web UI at http://localhost:5173
# You can:
# - View the list of available tools
# - Call tools directly with test inputs
# - Inspect the raw JSON-RPC messages
# - Verify tool schemas and responses
\end{lstlisting}

%-------------------------------------------------------------------------
\chapter{ANALYSIS AND REAL-WORLD IMPACT}
%-------------------------------------------------------------------------

\section{The Ecosystem Response}
The response to MCP's release in November 2024 was rapid and enthusiastic. Within weeks, major technology companies and open-source contributors had built MCP servers for a wide range of services. This ecosystem growth validated MCP's design: the protocol was simple enough to implement quickly but powerful enough to support complex use cases \cite{anthropic2024mcp}.

\begin{table}[H]
    \centering
    \caption{Notable MCP Servers in the Ecosystem (as of early 2025)}
    \label{tab:mcp_ecosystem}
    \begin{tabular}{|>{\raggedright\arraybackslash}p{3.0cm}|>{\raggedright\arraybackslash}p{3.5cm}|>{\raggedright\arraybackslash}p{6.0cm}|}
        \hline
        \textbf{Server} & \textbf{Provider} & \textbf{Capabilities Exposed} \\ \hline
        GitHub MCP & Anthropic (official) & Repository management, issues, pull requests, code search \\ \hline
        PostgreSQL MCP & Community & Database queries, schema inspection, data analysis \\ \hline
        Brave Search MCP & Brave & Web search, news search, local search \\ \hline
        Google Drive MCP & Community & File listing, reading, searching Google Drive \\ \hline
        Slack MCP & Community & Message sending, channel listing, user lookup \\ \hline
        Puppeteer MCP & Anthropic (official) & Web automation, screenshot capture, form filling \\ \hline
        Filesystem MCP & Anthropic (official) & Local file read/write, directory listing, search \\ \hline
        Fetch MCP & Anthropic (official) & HTTP requests, web page content retrieval \\ \hline
    \end{tabular}
\end{table}

\section{Adoption by Major Platforms}
The adoption of MCP by major development platforms accelerated its impact significantly. Microsoft integrated MCP support into VS Code's GitHub Copilot, making MCP tools available to millions of developers. Cursor, the AI-first code editor, adopted MCP as its primary tool integration mechanism. This platform adoption created a strong incentive for tool developers to build MCP servers: a single implementation could reach users across multiple major platforms.

\section{Comparison with Alternative Approaches}
MCP is not the only approach to AI tool integration, and it is worth comparing it to alternatives to understand its specific strengths and limitations.

\begin{table}[H]
    \centering
    \caption{MCP vs. Alternative AI Tool Integration Approaches}
    \label{tab:mcp_comparison}
    \begin{tabular}{|>{\raggedright\arraybackslash}p{2.5cm}|>{\raggedright\arraybackslash}p{2.5cm}|>{\raggedright\arraybackslash}p{2.5cm}|>{\raggedright\arraybackslash}p{2.5cm}|>{\raggedright\arraybackslash}p{2.5cm}|}
        \hline
        \textbf{Feature} & \textbf{MCP} & \textbf{OpenAI Function Calling} & \textbf{LangChain Tools} & \textbf{OpenAPI / Plugins} \\ \hline
        Vendor Neutral & Yes & No (OpenAI only) & Yes (library) & Partial \\ \hline
        Language Agnostic & Yes & Yes & No (Python) & Yes \\ \hline
        Standardized Protocol & Yes & No & No & Partial \\ \hline
        Ecosystem / Reuse & Growing rapidly & Limited & Large but fragmented & Limited \\ \hline
        Security Model & Built-in & Application-level & Application-level & Limited \\ \hline
        Remote Servers & Yes (SSE) & No & No & Yes \\ \hline
        Resource Access & Yes & No & Limited & No \\ \hline
        Prompt Templates & Yes & No & Partial & No \\ \hline
    \end{tabular}
\end{table}

\section{Limitations and Open Challenges}
Despite its strengths, MCP is not without limitations. Understanding these is important for a balanced assessment.

\textbf{Complexity for Simple Use Cases:} For a developer who only needs to connect one AI model to one tool, MCP introduces more structure than a simple function call. The protocol overhead (JSON-RPC, initialization handshake, capability discovery) is unnecessary for trivial integrations.

\textbf{Latency:} The MCP protocol adds a communication layer between the Host and the tool. For local stdio servers, this overhead is minimal. For remote SSE servers, the additional network round-trip can add meaningful latency to tool calls.

\textbf{Versioning and Evolution:} As MCP evolves, maintaining backward compatibility across protocol versions will be a challenge. The ecosystem's rapid growth means that many servers and clients will be built on early versions of the spec, and migration will require careful management.

\textbf{Discovery:} There is currently no universal registry for MCP servers. Users must manually configure which servers to connect to. A standardized discovery mechanism would significantly improve the user experience.

%-------------------------------------------------------------------------
\chapter{CONCLUSION}
%-------------------------------------------------------------------------

\section{Summary of Findings}
This report has traced the journey from the fragmented, chaotic world of pre-MCP AI tool integration to the standardized, composable ecosystem that MCP enables. The key findings are:

\textbf{The M×N problem was a genuine crisis.} The exponential growth of AI models and tools, combined with the absence of a standard integration protocol, was creating an unsustainable maintenance burden and limiting the practical capabilities of AI systems.

\textbf{MCP solves the right problem at the right level.} By defining a protocol rather than a library, MCP creates a stable, language-agnostic interface that decouples the development lifecycles of AI hosts and tool servers. The M×N problem becomes M+N.

\textbf{The three-layer architecture is well-designed.} The separation of Host, Client, and Server creates clear boundaries of responsibility, enables a strong security model, and makes the system composable and extensible.

\textbf{The three primitives cover the essential use cases.} Tools, Resources, and Prompts together address the full range of ways an AI assistant needs to interact with the external world: executing actions, reading data, and applying expertise.

\textbf{The ecosystem response validates the design.} The rapid adoption of MCP by major platforms and the growth of the server ecosystem demonstrate that the protocol addresses a real need in a practical way.

\section{Importance of the Topic}
MCP represents a foundational infrastructure shift in AI development. Just as TCP/IP enabled the internet by standardizing network communication, and just as USB enabled the personal computing ecosystem by standardizing peripheral connectivity, MCP has the potential to enable a new generation of AI applications by standardizing tool connectivity.

The significance of this cannot be overstated. An AI assistant that can reliably access any tool, any data source, and any service — through a single, standardized interface — is qualitatively more capable than one that is limited to its training data. MCP is the infrastructure that makes AI agents practical.

\section{Future Scope}
The MCP ecosystem is young and rapidly evolving. Several important directions for future development are evident:

\subsection{Standardized Server Discovery}
A universal registry or discovery mechanism for MCP servers would dramatically improve the user experience. Users should be able to search for and install MCP servers as easily as they install browser extensions or smartphone apps.

\subsection{Enhanced Security Primitives}
Future versions of MCP could include more sophisticated security primitives: fine-grained permission scopes, cryptographic authentication of servers, and standardized audit logging. These features would make MCP suitable for high-security enterprise environments.

\subsection{Multi-Agent Coordination}
As AI systems evolve toward multi-agent architectures — where multiple AI models collaborate on complex tasks — MCP could be extended to support agent-to-agent communication. An MCP server could expose the capabilities of another AI agent, enabling composable, hierarchical AI systems.

\subsection{Streaming and Long-Running Tools}
The current MCP specification handles request-response interactions well, but many real-world tools involve long-running operations (e.g., running a test suite, training a model, processing a large file). Future versions of MCP could standardize streaming progress updates and cancellation for long-running tool calls.

\subsection{Broader Language SDK Support}
While official SDKs exist for Python and TypeScript, the MCP ecosystem would benefit from official, well-maintained SDKs for other popular languages: Go, Rust, Java, and C\#. This would lower the barrier to building MCP servers for developers in these ecosystems.

\section{Closing Remarks}
The Model Context Protocol is more than a technical specification. It is a statement about how the AI ecosystem should be organized: open, interoperable, and composable. By providing a universal language for AI assistants to communicate with the world, MCP transforms AI from a powerful but isolated technology into a connected, capable participant in the digital infrastructure of modern life. The protocol is young, but its trajectory is clear: MCP is becoming the foundational layer on which the next generation of AI applications will be built.
